\section*{Appendix C: Ethics Statement and Broader Impacts}

\subsection*{C.1 Data Sources and Availability}

All datasets used in this work are publicly available:

\begin{itemize}
\item \textbf{Elliptic Bitcoin}~\cite{weber2019anti}: Publicly released for research. No personally identifiable information; transactions are anonymized.
\item \textbf{Amazon Fraud}~\cite{mcauley2015amazon}: Public Amazon product co-purchase data with synthetic fraud labels.
\item \textbf{YelpChi}~\cite{rayana2015yelp}: Publicly available Yelp review dataset with spam labels from Yelp's filter.
\item \textbf{DGraphFin}~\cite{huang2022dgraph}: Public benchmark released through PyTorch Geometric. Features are anonymized.
\item \textbf{Citation networks} (Cora, Citeseer, Pubmed): Standard public benchmarks.
\item \textbf{Heterophily benchmarks} (Texas, Wisconsin, Cornell): Standard public benchmarks.
\end{itemize}

\textbf{No private data}: This work does not use any proprietary or non-public data. All experiments can be reproduced using publicly available datasets.

\subsection*{C.2 Ethical Considerations}

\textbf{Potential benefits}:
\begin{itemize}
\item Improved fraud detection protects consumers and financial institutions
\item The FSD framework provides interpretable guidance for model selection, reducing arbitrary choices
\item Transparent $\rho_{\text{FS}}$ metric enables practitioners to understand \textit{why} certain methods work
\end{itemize}

\textbf{Potential risks}:
\begin{itemize}
\item \textbf{False positives}: Fraud detection models may incorrectly flag legitimate users, causing inconvenience or financial harm
\item \textbf{Adversarial adaptation}: Fraudsters may adapt behavior to evade detection; regular model retraining is essential
\item \textbf{Dual-use}: Techniques for detecting fraud could theoretically be reversed to improve evasion strategies
\end{itemize}

\textbf{Mitigation strategies}:
\begin{itemize}
\item Model predictions should supplement, not replace, human expert review
\item Regular fairness audits across user demographics
\item Provide recourse mechanisms for users incorrectly flagged as fraudulent
\item Monitor model performance drift over time
\end{itemize}

\subsection*{C.3 Reproducibility Statement}

To ensure reproducibility:

\begin{enumerate}
\item \textbf{Code}: Complete implementation released under MIT license at \texttt{https://github.com/MengdanXue/trust-regions-gnn}
\item \textbf{Hyperparameters}: All training details provided in Appendix B
\item \textbf{Random seeds}: 5-seed experiments use seeds 42, 123, 456, 789, 1024
\item \textbf{Statistical tests}: Paired t-tests with p-values reported for all comparisons
\item \textbf{Hardware}: NVIDIA RTX 3090 (24GB VRAM), PyTorch 2.0, PyG 2.3, CUDA 11.8
\end{enumerate}

\subsection*{C.4 Computational Resources}

\textbf{Training costs}:
\begin{itemize}
\item Single model training: 1-10 minutes per dataset (RTX 3090)
\item Full experiment suite (all models, all datasets, 5 seeds): $\sim$8 hours
\item $\rho_{\text{FS}}$ computation: 2-5 minutes per dataset
\item Estimated carbon footprint: $<$5 kg CO$_2$ equivalent (minimal compared to large-scale pretraining)
\end{itemize}

\subsection*{C.5 Limitations}

We acknowledge the following limitations:

\begin{enumerate}
\item \textbf{Dataset coverage}: Four financial fraud datasets may not cover all real-world fraud scenarios
\item \textbf{Simplified theory}: Theorem~\ref{thm:optimal_regime} provides qualitative guidance but not precise thresholds
\item \textbf{Static graphs}: Our analysis assumes static graphs; extension to dynamic/temporal graphs requires additional work
\item \textbf{Scalability}: NAA's 3$\times$ computational overhead may be prohibitive for billion-scale graphs
\item \textbf{Feature requirements}: NAA assumes meaningful numerical features; may not apply to categorical-only domains
\end{enumerate}

\subsection*{C.6 Author Contributions}

[REDACTED FOR DOUBLE-BLIND REVIEW]

\subsection*{C.7 Conflicts of Interest}

The authors declare no financial or commercial relationships that could constitute a conflict of interest.
